\documentclass[12pt,conference,onecolumn]{IEEEtran}

\usepackage[hidelinks]{hyperref}

\title{Optimizing the fabrication process of microspheres for extended-release drug delivery: internship at ACON Pharmaceuticals}
\author{%
\IEEEauthorblockN{Emily Chen}\IEEEauthorblockA{Science \& Engineering\\Manalapan High School\\Englishtown, NJ\\\href{mailto:426echen@frhsd.com}{426echen@frhsd.com}}\and
\IEEEauthorblockN{Petra Rofman}\IEEEauthorblockA{Science \& Engineering\\Manalapan High School\\Englishtown, NJ\\\href{mailto:426profman@frhsd.com}{426profman@frhsd.com}}
}
\date{June 16, 2026}

\usepackage{siunitx}
\newcommand{\keywords}{microspheres, emulsion-evaporation, co-acervation, encapsulation rates, the China Method, solidification, fracture}

\usepackage{hyperref}
\makeatletter
\AtBeginDocument{
\hypersetup{%
pdftitle={\@title},
pdfauthor={Emily Chen and Petra Rofman},
pdfkeywords={\keywords}}}
\makeatother

\begin{document}
\maketitle 

\begin{abstract}
Frequently and consistently taking pills to treat schizophrenia can be inconvenient or impossible for many patients. As a more reliable method to ensure patients consistently get their medicine, ACON Pharmaceuticals works to develop a procedure to make microspheres, small quantities of drug encased within polymers, for long term diffusion of medicine into the bloodstream. We are conducting drug dissolution studies with the microspheres in order to achieve a drug release over the span of one month. With this internship we will also explore two different microsphere-making procedures. 

The first one focuses on a process called ``The China Method.'' Currently, on a 10x scale, we have maximized the percentage of drugs loaded into the microspheres. This process employs the optimal drug to polymer ratio and most effective stirring methods to break down or dissolve the polymers. These details were determined through trial and error experiments. Our goal is to upscale this method and adjust it accordingly so that it will continue to maximize drug loading on a larger scale, thus being ideal for mass production. 
The second method is a new procedure that we will explore with the purpose of finding a more time or material efficient technique to streamline the mass production of microspheres. This method requires melting the polymer and drug together, allowing the mixture to cool into a solid, and then breaking the solid apart with a rotational milling machine. Our goal is to experiment with the ratio of drug to polymer as well as the length and speed of the milling process to optimize this new method.

%In the pharmaceutical industry, numerous drugs are required to be released over extended periods of time for diverse applications. Microspheres, microparticles that contain drugs encapsulated by polymers, have special characteristics that make them optimal for these applications. The quality of microspheres is characterized by their physical shape/size and drug encapsulation rates. My goal is to optimize the fabrication process's efficiency and improve the quality of microspheres produced by these techniques by producing microspheres that have an average diameter of about \qtyrange{10}{20}{\micro\meter} and have as high of a drug encapsulation rate as possible.

%In the first semester of my internship at ACON Pharmaceuticals, I primarily focused on testing the emulsion-evaporation fabrication technique. Using this time-consuming method, which involves the evaporation of an emulsified drug-polymer solution over 2-3 days, I found that the microspheres were suboptimal, often consisting of excess polymer and low drug encapsulation rates.

%For my second semester, I tested using the phase-separation method of fabrication, where I rapidly induced the formation of microspheres within minutes by adding a liquid to the emulsified drug-polymer solution. This method proved to be time-efficient without sacrificing the quality of the microspheres. To optimize for mass production, I aimed to enhance the drug encapsulation efficiency of this process by testing adding a base, and I also aimed to minimize the amount of materials used, while also scaling up and making larger batches. Furthermore, I began drug dissolution studies on the produced microspheres to measure the extended-release rates for drug delivery.
\end{abstract}

\begin{IEEEkeywords}
\keywords
\end{IEEEkeywords}

\end{document}